% Nouveau EGA V (cote 25): archive cote 25, site page 88 (PDF page 89).
% typescript. Draft compared: none.
% AI-transcribed by Claude Opus 5.5 (claude-opus-5-5, adaptive thinking, effort high), Message Batches API, batch ega25b (msgbatch_01A7n5YuXubUPNK1zsepi8RG), 2026-11-13; checked against both supplied scan parts; see CHECK_DECISIONS.tsv.
\scanpage{25/88}
\mspage{\hand{82}}

Grass$^{m}$(P), la dimension de $\underline{X}^{(m)}_{\hand{\xi}}$ est au moins égale à dim X-m , il en est
% NOTE: the right edge of the scan cuts the ends of the first two lines ("est", "l'inéga")
en particulier ainsi lorsque $\hand{\xi}$ provient d'un point de Q(C$_{\hand{\unsure{\eta|y}}}$). Pour l'inéga
lité en sens inverse, notons que (nous plaçant sur le corps k' = k($\hand{\xi}$))
comme $C_{\hand{\unsure{\eta|y}}\,k'} \hand{\subset} L_{\hand{\xi}}$ est  un hyperplan de $L_{\hand{\xi}}$ , si la dimension de
$\underline{X}^{(m)}_{\hand{\xi}} = Xx_{P}L_{\hand{\xi}}$ était \hand{$\geqslant$} dim X - m + 1 , alors celle de $\underline{X}^{(m+1)}_{\hand{\unsure{\eta|y}}}$ =
$Xx_{P}C_{\hand{\unsure{\eta|y}}}$ serait \hand{$\geqslant$} dim X - m , (puisque le changement de base k($\hand{\unsure{\eta|y}}$)$\hand{\to}$k'
transforme ce dernier préschéma en $(Xx_{P}L_{\hand{\xi}})x_{L_{\hand{\xi}}}(C_{\hand{\unsure{\eta|y}}\,k'})$ ). Or on a au
contraire  dim $\underline{X}^{(m+1)}$ = dim X - m -1 
, en vertu du n°2 (revu au n°10).
Le cas b)  se traite
de façon analogue: si on avait dim $Xx_{P}L$ \hand{$\geqslant$} d\hand{+1},  ou
ce qui revient au même  $f_{k'}(X_{k'}) \hand{\cap} L_{\hand{\xi}}$ de dimension \hand{$\geqslant$}1,
alors on aurait par le même argument que plus  haut que
$\underline{X}^{(m+1)}_{\hand{\unsure{\eta|y}}} \hand{\neq} \hand{\emptyset}$ , contrairement à ce qu'on avait signalé avant 14.1.

\uline{Corollaire} 14.2. Supposons X de dimension m et f:X$\hand{\to}$P fini (resp quasi-fini), alors
le morphisme $p_{C_{\hand{\unsure{\eta|y}}}} : X_{k(\hand{\unsure{\eta|y}})} \hand{\unsure{\longrightarrow|\relbar}} Q(C_{\hand{\unsure{\eta|y}}})$ est fini surjectif(resp. quasi-fini
dominant).

En effet, ce morphisme est quasi-fini, et comme dim $X_{\hand{k(\unsure{\eta|y})}}$ = dim Q (C$_{\hand{\unsure{\eta|y}}}$), il est
dominant;  si f est fini, $p_{C_{\hand{\unsure{\eta|y}}}}$ est de plus fini,
 car propre, donc surjectif étant dominant.
\uline{Corollaire} 14.3. Sous les conditions de 14.1. a), si X est Cohen-Macau-%
lay, alors le morphisme $p_{C} : \overset{\hand{\sim}}{X}(C_{\hand{\unsure{\eta|y}}}) \hand{\unsure{\longrightarrow|\relbar}} Q(C_{\hand{\unsure{\eta|y}}})$ est un morphisme de
Cohen-Macaulay, et à fortiori \uline{plat}.

Pour la démonstration, cf remarque en haut de la page 21, avant 5.6. qui donne un résultat plus fort, (à inclure dans 14.3), savoir que les $\varphi^{(m)}_{\xi}$ pour $\xi\in Q(C_{\unsure{\eta|y}})$ sont \unsure{$F_{k(\xi)}$|$F_{\xi}$}-réguliers.
Ce Corollaire devrait être modulé, mais pour simplifier on pourra
supposer f quasi-fini: si F est un Module de Cohen-Macaulay sur X, et si
pour to\uline{u}te composante irréductible Z de Supp F, on a dim Z \hand{$\geqslant$} m, alors
$\overset{\hand{\sim}}{F}(C_{\hand{\unsure{\eta|y}}})$ est CM (et à fortiori plat) relativement à Q(C$_{\hand{\unsure{\eta|y}}}$).

On notera qu'on ne peut pas remplacer, pour obtenir la

\placedmargin{88}
